\section{Introduction}\label{sec:introduction}

A Ricci flow is a family of Riemannian metrics satisfying
\(\partial_tg=-2\Ric(g)\). Its scalar curvature \(R=\tr_g\Ric\)
controls the infinitesimal change of volume, whereas its full curvature
tensor controls smooth continuation. The scalar-curvature extension
problem asks whether a uniform scalar bound prevents a finite-time
singularity. We prove its positive resolution on closed manifolds in
real dimension four.

\begin{theorem}[Smooth extension]\label{thm:extension}
Let \(M\) be a closed connected smooth four-manifold, let
\(0<T<\infty\), and let \(g(t)\), \(0\le t<T\), be a smooth Ricci
flow with smooth initial metric. If
\[
 \sup_{M\times[0,T)}|R(g(t))|<\infty,
\]
then there are \(\varepsilon>0\) and a smooth Ricci flow
\(\widetilde g(t)\) on \(M\), \(0\le t<T+\varepsilon\), such that
\(\widetilde g(t)=g(t)\) for every \(t<T\).
\end{theorem}

The same conclusion holds without connectedness. The closed manifold
\(M\) has finitely many connected components, each a closed smooth
four-manifold. The restricted flows satisfy the hypotheses of
Theorem~\ref{thm:extension} at the common time \(T\), so each extends
to \(T+\varepsilon_j\) for some \(\varepsilon_j>0\). Taking
\(\varepsilon=\min_j\varepsilon_j\) and combining these flows gives
the required extension on \(M\).

\begin{corollary}\label{cor:scalar-blowup}
If a Ricci flow on a closed smooth four-manifold has finite maximal
smooth existence time \(T\), then
\[
 \limsup_{t\uparrow T}\max_{x\in M}R(g(t))(x)=+\infty.
\]
\end{corollary}

\begin{proof}
The equation \((\partial_t-\Delta)R=2|\Ric|^2\) and the maximum
principle give a scalar lower bound on a finite time interval. A finite
upper limit in the assertion would therefore give a uniform absolute
bound near \(T\), and smoothness supplies the bound on the remaining
compact time interval. The componentwise extension just described
contradicts maximality.
\end{proof}

\subsection{The extension problem and its geometric obstruction}

Hamilton introduced Ricci flow and established its short-time existence
on closed manifolds \cite{Hamilton1982}. A compact flow continues while
its full curvature remains bounded. \v{S}e\v{s}um showed that a uniform
bound for the Ricci tensor also suffices \cite{Sesum2005}. A scalar bound
is weaker: in a possible curvature blowup, the scalar curvature may
vanish after rescaling while the full curvature stays nonzero. The
limiting geometry can therefore be Ricci-flat without being flat.

Bamler and Zhang developed heat-kernel and curvature estimates under
scalar bounds, including the four-dimensional energy and compactness
theory \cite{BamlerZhangI,BamlerZhangII}. Their work also gives smooth
extension when \(H_2(M;\mathbb F)=0\) for every field \(\mathbb F\)
\cite[arXiv version, Corollary~1.10]{BamlerZhangI}. Simon constructed
continuation from the limiting space by an orbifold Ricci flow
\cite{SimonExtension}. Bamler's convergence theory gives smooth
convergence away from a set of codimension at least four
\cite{BamlerConvergence};
four-dimensional convergence under spatial scalar-integrability
hypotheses is also studied by Liu and Simon \cite{LiuSimon}.
These compactness conclusions permit concentrated curvature and
orbifold points. The assertion here is smooth continuation of the given
flow on its original manifold. The scalar extension problem and its
known special cases are discussed in \cite[Section~2]{LiLiSurvey}.
Finite-time scalar blowup was established for closed K\"ahler--Ricci
flows by Zhang \cite{Zhang2010Kahler} and for compact Ricci flows with
a global Type~I bound, $\sup_M|\Rm(g(t))|\le C/(T-t)$, by
Enders, M\"uller and Topping \cite{EndersMuellerTopping2011}.
More recently, Buzano and Di Matteo proved scalar blowup at local
Type~I singular points without a global Type~I assumption
\cite[Definition~1.4 and Theorem~1.1]{BuzanoDiMatteo2026}.
Here ``local Type~I'' means
\(0<\limsup_{t\uparrow T}(T-t)r_{\Rm}(p,t)^{-2}<\infty\),
where \(r_{\Rm}\) is their parabolic curvature scale at the point \(p\).
That result addresses a specified blowup regime; the argument below
allows unrestricted rates at the hypothetical concentration points.

The spatial degeneration analysis has its origins in the theory of
Einstein metrics. Anderson's compactness theorem produces orbifold
limits with smooth convergence on their regular parts
\cite[Theorem~C]{Anderson1989}. Bando's bubbling analysis, including
its correction, organizes Ricci-flat bubbles and quantitative necks
\cite{Bando1990,BandoCorrection1990}; the work of Bando, Kasue and
Nakajima develops coordinates at infinity under curvature decay and
volume-growth hypotheses \cite{BandoKasueNakajima1989}.
Cheeger and Tian's curvature-energy estimates provide a further
important part of this Einstein four-manifold theory
\cite{CheegerTian2006}. Our slices need not be Einstein. We prove
the required geometric statements below using the scalar-bound
estimates and retaining the Ricci forcing in the neck argument.

At a hypothetical concentration point, successive spatial rescalings
produce a finite tree of Ricci-flat pieces. A vertex represents a
nonflat limit; its singular points may contain smaller vertices. The
outermost nonflat vertex has a distinguished length scale \(a(t)\).
The central problem is to rule out contraction of this outermost scale
while allowing arbitrarily smaller descendants.

\subsection{The static functional and the flow estimates}

Retain the true metric from the cluster through an annular collar of
radius \(b=Na\), with \(N\) large and fixed, and attach an
asymptotically locally Euclidean end. On that end write the metric in
Euclidean quotient coordinates. The renormalized Einstein--Hilbert
functional is
\[
 E(g)=\lim_{D\to\infty}\left[
 \int_{\{r<D\}}R(g)\,d\mu_g-
 \int_{\{r=D\}}(\partial_jg_{ij}-\partial_ig_{jj})n_i\,dS_\delta
 \right].
\]
The boundary subtraction cancels the scalar curvature's linear
divergence term. The functional scales by the square of the length
unit, and its first variation pairs the metric velocity with
\(\tfrac12 Rg-\Ric\).

The companion article \cite{StaticTrees} proves a sequential inequality
\(E=o(M)\) on each fixed limiting tree, where \(M\) is its weighted,
dimensionless Ricci error. Section~\ref{sec:static-input} records the
complete input, including the end weights, derivative orders and
matching quotient groups. Our task is both to verify those hypotheses
for metrics obtained from the flow and to compare their functional
variation with actual spacetime Ricci energy.

\begin{figure}[htbp]
 \centering
 \begin{tikzpicture}[
  x=1cm,y=1cm,>=Stealth,
  piece/.style={draw=blue!55!black,fill=blue!5,rounded corners=3pt,
                minimum width=2.1cm,minimum height=.65cm,align=center},
  edge/.style={draw=blue!55!black,thick},
  every node/.style={font=\small}]
  \node[piece] (root) at (0,0) {outermost piece\\scale $a$};
  \node[piece] (left) at (-2.6,-1.65) {descendant\\$a e^{-2L_1}$};
  \node[piece] (right) at (2.6,-1.65) {descendant\\$a e^{-2L_2}$};
  \node[piece] (deep) at (-2.6,-3.2)
    {further descendant\\$a e^{-2L_1-2L_3}$};
  \draw[edge,->] (root) -- node[above left] {$L_1$} (left);
  \draw[edge,->] (root) -- node[above right] {$L_2$} (right);
  \draw[edge,->] (left) -- node[right] {$L_3$} (deep);
  \draw[edge] (root.north) -- (0,.9);
  \draw[draw=orange!75!black,thick] (-1.5,.9)--(1.5,.9);
  \node[right,align=left,text=orange!65!black] at (1.65,.9)
    {fixed collar\\$b=Na$};
  \draw[edge,dashed,->] (0,.9)--(0,1.7);
  \node[above] at (0,1.7) {attached ALE end};
\end{tikzpicture}

 \caption{The outermost scale \(a\), its smaller descendants, and the
 material collar at radius \(b=Na\). An edge of logarithmic length
 \(L\) changes the length unit by \(e^{-2L}\). The collar extension
 preserves the inner geometry and supplies an ALE end. Incidence and
 scale are schematic.}
 \label{fig:scales}
\end{figure}

Two flow estimates are decisive. The divergence constraint on Ricci
removes the critical radial harmonic mode on a nearly flat annulus.
This improves annular decay enough that the loss outside the collar
is small compared with the energy inside the root. The exterior is
then constructed as an analytic function of the metric on one fixed
material collar. Longer exterior coordinates enter only an error
estimate. Consequently the extension can be differentiated on time
patches without choosing a smoothly varying bubble tree.
The collar construction uses the classical compatibility between
metric strains and linearized curvature, together with Korn control
modulo rigid motions \cite{Khavkine2017,CiarletCiarlet2005}.
The scale-specific estimates, filled-space inverse and dependence
on the fixed collar are proved in Section~\ref{sec:collar}.

\subsection{Proof structure}

Section~\ref{sec:geometry} establishes the scale estimates, genuine
orbifold limits and finite trees needed here. It defines a continuous
outermost scale and selects buffered time intervals on which that
scale changes by a fixed factor. After parabolic normalization, the
interval has fixed duration and its root length \(a\) tends to zero.

Section~\ref{sec:parabolic} controls Ricci on exterior annuli and on
the root's compact and joining regions. A time maximal function
converts root spacetime energy into the single-time weighted error
required by the static theorem. Section~\ref{sec:collar} constructs
the collar-dependent exterior and proves an error bound independent
of its length.

Section~\ref{sec:contradiction} applies the static theorem separately
to arbitrary sequences of active time patches. If \(K^2\) denotes
the root spacetime Ricci energy divided by \(a^4\), the positive
functional variation and its static smallness imply
\[
 cK^2\le o(1)(K+1)^2.
\]
Hence \(K\to0\), without assuming it was bounded. Dense sets of
material points then preserve distances at scale \(a\) between the
two endpoints. The resulting pointed isometry preserves the
curvature-radius maximum defining the outermost scale. This contradicts
its prescribed factor change. The absence of concentration gives a
smooth terminal metric and continuation on \(M\).

\subsection{Conventions}

We use \(\Delta=\tr\nabla^2\). A tensor norm at scale \(r\)
includes \(r^j\) on its \(j\)th spatial derivative, while leaving
its pointwise magnitude unchanged. Tensor magnitudes use the current
metric unless specified otherwise. Quotient-chart components are
Cartesian components on the Euclidean cover; all constructions are
equivariant and descend to the quotient.

Constants may depend on a fixed finite differentiation order. After
extraction they may also depend on a fixed finite tree. Uniformity
in a sequence index, a length, or the collar factor \(N\) is stated
where needed. A sequential little-oh is taken after fixing the
parameters in its statement.
